\subsection{Векторные диаграммы}

%%%%%%%%%%%%%%%%%%%%%%%%%%%%%%%%%%%%%%%%%%%%%%%%%%%%
%%%%%%%%%%%%%%%%%%%%%%%%%%%%%%%%%%%%%%%%%%%%%%%%%%%%
 
 \z{ Со скалы, возвышающейся над морем на высоту $h = 25\, \text{м}$, бросили камень. Найдите время его полёта, если известно, что непосредственно перед падением в воду камень имел скорость $v = 30\, \text{м/с}$, направленную под углом $\beta = 120\degree$ к начальной скорости.}
%
{$t=4.4\,\text{с}$}

%%%%%%%%%%%%%%%%%%%%%%%%%%%%%%%%%%%%%%%%%%%%%%%%%%%%

 \z{Камень бросили под углом к горизонту с начальной скоростью $v_0  = 25\, \text{м/с}$. Через время $\tau$ он достиг максимальной высоты, удалившись по горизонтали на расстояние $L = 30\, \text{м}$ от места броска. Найдите время $\tau$.  }
%
{$t=\sqrt{\frac{v_0^2\pm\sqrt{v_0^4-4g^2L^2}}{2g^2}}$}
	
%%%%%%%%%%%%%%%%%%%%%%%%%%%%%%%%%%%%%%%%%%%%%%%%%%%%
\z{Скорость камня $v_0$, брошенного под углом $\varphi = 60\degree$ к горизонту, уменьшилась вдвое за $\Delta t = 1\, \text{с}$. Найдите модуль перемещения $S$, которое за это время совершил камень.}
%
{$S=g\Delta t^2\sqrt{\frac{7}{12}}\approx 7.62\,\text{м}$}

%%%%%%%%%%%%%%%%%%%%%%%%%%%%%%%%%%%%%%%%%%%%%%%%%%%%	
\z[{\begin{wrapfigure}[5]{r}{0.3\textwidth}%
		\vspace{-20pt}%
		\begin{flushright}%
 			\includegraphics[width=0.3\textwidth]{2.6.4}%
		\end{flushright}%
	\end{wrapfigure}%
	}]%
	{Кот Леопольд сидел у края крыши. Два злобных мышонка выстрелили в него из рогатки. Однако камень, описав дугу, упал у ног кота (см. рисунок) через время $\tau = 1\,\text{с}$. На каком расстоянии $s$ от мышей находился кот Леопольд, если известно, что векторы скоростей камня в момент выстрела и в момент падения были взаимно перпендикулярны? }
%
{$s=\frac{g\tau^2}{2}$}

%%%%%%%%%%%%%%%%%%%%%%%%%%%%%%%%%%%%%%%%%%%%%%%%%%%%	
\z{Мальчик бросил камень под некоторым углом $\alpha$ к горизонту. Пренебрегая сопротивлением воздуха, определите, при каких значениях угла бросания $\alpha$ камень всё время (до падения на землю) будет удаляться от мальчика. }
%
{$\sin{\alpha}<\frac{2\sqrt{2}}{3}$}
